% Source: https://artofproblemsolving.com/community/c3h2501903_some_problems
% Stable identity: aops-topic-2501903-q07
\textbf{Yes.} Take
\[P(x)=x^2+x+4.\]
Its coefficients $1,1,4$ have no common factor greater than $1$. For every integer $n$, the number $n^2+n=n(n+1)$ is a product of two consecutive integers, so it is even and nonnegative. Hence $P(n)=n(n+1)+4$ is even and at least $4$, so $|P(n)|=P(n)$ is composite.

(Examples exist in every degree $d\ge2$. For even $d$, $x^d+x+4$ works: $n^d+n$ is even and $n^d\ge|n|$, so every value is even and at least $4$. For odd $d$, $x^d-x+12$ works: $n^d\equiv n$ modulo $2$ and modulo $3$, so $6\mid n^d-n$, and $n^d-n\ne-12$ (it is $0$ for $|n|\le1$, equals $\pm\left(2^d-2\right)\ne\pm12$ for $|n|=2$, and has absolute value at least $3\left(3^2-1\right)=24$ for $|n|\ge3$), so every value is a nonzero multiple of $6$.) $\blacksquare$
